\documentclass[12pt]{article}
\input{.house-style/preamble.tex}
\usepackage{xurl}
\usepackage{float}
\setlength{\headheight}{14pt}
\clubpenalty=10000
\widowpenalty=10000
\emergencystretch=1em
\AtBeginBibliography{\emergencystretch=1em}
\setcounter{biburlnumpenalty}{100}
\newcommand{\tablebody}[1]{\csname @@input\endcsname #1 }
\providecommand{\censusClaims}{658}
\providecommand{\withdrawnClaims}{7}
\providecommand{\liveClaims}{651}
\providecommand{\taxonomicClaims}{130}
\providecommand{\otherClaims}{6}
\providecommand{\lexemelessOnly}{79}
\providecommand{\lexemelessOverlap}{98}
\providecommand{\retainedClaims}{436}
\providecommand{\claimPairs}{450}
\providecommand{\keyedClaims}{51}
\providecommand{\keyedAgree}{45}
\providecommand{\labelModel}{Claude Haiku 4.5}
\providecommand{\typesafeAccepted}{252}
\providecommand{\typesafeAgree}{237}

\hypersetup{pdftitle={Quantificational controls: attestations and provenance}}
\title{Claim register: participation claims by diagnostic, lexeme, status, and basis}
\author{Brett Reynolds}
\date{Supplementary material, September 2026}
\fancyhead[L]{\small\scshape Claim register}
\usepackage{longtable}
\begin{document}
\maketitle

\section{What this register is}

This supplement lists every participation claim the article makes, by diagnostic and lexeme, with the claim's status (licensed, conditional, excluded) and the basis the article gives for it at that point in the text. It exists so that the comparison in §2 of \textit{Determinatives as nouns in English} can be inspected: for each diagnostic, which lexemes the article treats as showing the property, which as excluded, and on what evidence. It is a register of what the article asserts, not an independent survey of English, and it doesn't assess whether any claim is correct.

\section{How it was made}

A census of the article's sentences was produced on 14 September 2026 by a language model (Aristotle, Harmonic) reading each section and recording every claim that an expression participates in a construction, with an exact quotation: \censusClaims{} claims across all 28 sections, every quotation an exact substring of the source, and ten withheld acceptance cases all reached. The census is model-identified: the acceptance cases and the keyed claims below test recall of known claims and agreement of labels, not the rate of omitted or spurious claims, which no human audit has estimated. Each claim was then labelled with a diagnostic from a 22-item catalogue by \labelModel{} in two passes; against \keyedClaims{} claims with an independent key, the label was reproduced for \keyedAgree{}. A second model (TypeSafe jev-1.13.0) labelled all \censusClaims{} claims independently on 16 September 2026; on the \typesafeAccepted{} claims where it reported confidence of 0.8 or above, it agrees with the current labels on \typesafeAgree{}. The basis of each claim is the article's own: an attestation footnote, a starred form, or a citation at the sentence, and otherwise the census's reading of the sentence. Since the census, \withdrawnClaims{} claims have been withdrawn with the sentences they quoted (listed with the reasons in \texttt{amendments.json}), leaving \liveClaims{}. Of these, \taxonomicClaims{} are classed as taxonomic or methodological rather than constructional, \otherClaims{} are labelled only \enquote{other}, and a further \lexemelessOnly{} have no resolvable lexeme (\lexemelessOverlap{} more lexeme-less claims fall in the first two classes); the register covers the remaining \retainedClaims{} claims as \claimPairs{} claim--lexeme pairs, since a claim that names two lexemes appears once under each. Every quotation is re-checked against the current text of the article at each build, so the register can't drift from the prose it describes. Files, scripts, and model outputs are in the project repository under \texttt{analysis/manuscript-census-2026-09-14/}.

The labels are model-assigned and the key covers a twelfth of the claims, so a row is a pointer to the sentences it counts, not a verdict. The sections column gives the article sections to read.

\section{Coverage by diagnostic}

Table~\ref{tab:claim-coverage} gives, for each diagnostic, the number of lexemes of each subcategory with a claim, and in parentheses how many of them are claimed licensed or conditional / excluded only.

\begin{table}[H]
\centering\footnotesize
\caption{Lexemes with a claim, by diagnostic and subcategory; in parentheses, licensed or conditional / excluded only. The first two columns state whether the property is general, restricted, or absent for English common nouns and for adjectives (fused-head uses and the article's quantificational-adjective controls included), each cell sourced to \textit{CGEL} or to the article in \texttt{analysis/typicality.json}; a question mark marks a cell for which no source was located. Blocks follow the article's argument (§2.1, §§2.2 and 2.5, §2.4, §4); within a block, rows run from the properties that most favour Noun (general for nouns, absent for adjectives) to those that most favour Adjective, with unsourced rows last. Subcategory columns follow Table~2 of the article, with the adjective controls, which lie outside Noun, last.}\label{tab:claim-coverage}
\setlength{\tabcolsep}{3pt}
\begin{tabular}{>{\raggedright\arraybackslash}p{3.0cm}*{2}{>{\centering\arraybackslash}p{1.25cm}}*{5}{>{\centering\arraybackslash}p{1.35cm}}>{\centering\arraybackslash}p{0.7cm}}
\toprule
Diagnostic & of common nouns & of adjectives & common noun & proper noun & pronoun & determ. & adjective (control) & total \\
\midrule
\tablebody{analysis/generated/claim-coverage.tex}
\bottomrule
\end{tabular}
\end{table}

\section{The register}

\small
\setlength{\tabcolsep}{3pt}
\begin{longtable}{>{\raggedright\arraybackslash}p{2.7cm}>{\raggedright\arraybackslash}p{1.5cm}>{\raggedright\arraybackslash}p{2.15cm}>{\raggedright\arraybackslash}p{2.0cm}>{\raggedright\arraybackslash}p{2.6cm}>{\raggedright\arraybackslash}p{2.2cm}r}
\caption{Participation claims by diagnostic and lexeme, in the order of Table~\ref{tab:claim-coverage}; within a diagnostic, by subcategory and then alphabetically.}\label{tab:claim-register}\\
\toprule
Diagnostic & Lexeme & Subcategory & Status & Basis & Sections & Claims \\
\midrule
\endfirsthead
\toprule
Diagnostic & Lexeme & Subcategory & Status & Basis & Sections & Claims \\
\midrule
\endhead
\bottomrule
\endfoot
\tablebody{analysis/generated/claim-register.tex}
\end{longtable}

\printbibliography
\end{document}
